\section{The energy and the construction}

Write $\lambda(n)=(-1)^{\Omega(n)}$. For a sign function $f$ and a real
$x\ge0$, put $L_f(x)=\sum_{n\le x}f(n)$, with $L_f(x)=0$ for $x<1$.
For integers $N\ge1$, its signed Fej\'er energy is
\[
 Q_f(N)=\frac1N\sum_{t\in\mathbb Z}
     \left|\sum_{\substack{t<n\le t+N\\1\le n\le N}} f(n)\right|^2.
\]
This is a finite sum. A pair of indices at distance $h$ belongs to
exactly $N-h$ windows. Thus the expansion weights the signed pairs
at that distance by $2(1-h/N)$.

\begin{lemma}\label{round26:lem:prefix}
For every real sign function and every $N\ge1$,
\begin{align}
 Q_f(N)&=\frac{|L_f(N)|^2}{N}
  +\frac1N\sum_{j=1}^{N-1}
     \left(|L_f(j)|^2+|L_f(N)-L_f(j)|^2\right),\label{round26:eq:prefix}\\
 \frac12|L_f(N)|^2&\le Q_f(N)
      \le5\max_{0\le j\le N}|L_f(j)|^2.\label{round26:eq:prefix-ineq}
\end{align}
\end{lemma}
\begin{proof}
The window beginning at $t=0$ contains the full truncation. Each
nonempty proper prefix occurs once at $t=j-N$, and its complementary
suffix occurs once at $t=j$. All other windows have zero sum.
This gives \eqref{round26:eq:prefix}. The identity
\[
 |u|^2+|v-u|^2=2|u-v/2|^2+|v|^2/2
\]
gives the lower bound. If $M=\max_{j\le N}|L_f(j)|$, the summand in
\eqref{round26:eq:prefix} is at most $M^2+4M^2$, which gives the upper bound.
\end{proof}

Let $\Theta$ be the supremum of the real parts of the nontrivial
zeros of $\zeta$. For a fixed cutoff $y_0\ge2$, put
\[
 c_{y_0}(s)=\prod_{p\le y_0}\frac{1-p^{-s}}{1+p^{-s}}.
\]
This is a nonvanishing holomorphic function on $\R s>0$.
For the Dirichlet series $F_f$ of $f$, write
$H_f(s)=F_f(s)/(\zeta(s)c_{y_0}(s))$.

\begin{theorem}\label{round26:thm:construction}
For every $\beta\in(1/2,1)$ and $y_0\ge2$, there is one completely
multiplicative $f:\mathbb N\to\{-1,1\}$ with the following properties.
\begin{enumerate}
\item $f(p)=\lambda(p)$ for $p\le y_0$ and
  $\sum_{p:f(p)=1}p^{-1}<\infty$.
\item For every $A>0$, there is $C_{A,f}$ such that
  $|L_f(x)|\le C_{A,f}x/(\log x)^A$ for all $x\ge2$.
\item The Dirichlet series $F_f(s)=\sum_{n\ge1}f(n)n^{-s}$ extends
  meromorphically to $\R s>1/2$ as
  \begin{equation}\label{round26:eq:factorization}
   F_f(s)=\frac{\zeta(2s)\zeta(s+1-\beta)}{\zeta(s)}R_f(s),
  \end{equation}
  where $R_f$ is holomorphic and nonvanishing there and positive on the
  real axis. The function $F_f$ has a simple zero at~$1$, and $H_f$
  has a double zero there. Both have a simple pole at~$\beta$.
\item The energy satisfies
  \begin{equation}\label{round26:eq:exponent}
   \limsup_{N\to\infty}\frac{\log(1+Q_f(N))}{\log N}
      =2\max(\Theta,\beta).
  \end{equation}
  In particular, $Q_f(N)/N^{1+\varepsilon}$ is unbounded for every
  $0<\varepsilon<2\beta-1$.
\end{enumerate}
\end{theorem}

The same $f$ works for all $A$ and $\varepsilon$ in the theorem.\footnote{The construction is one fixed completely multiplicative function after choosing $\beta$ and the initial prime cutoff. It is not a different comparison sequence at each value of $N$.}
The construction in the next section rounds prime weights by taking
integer parts of their cumulative sum. The proof uses classical
logarithmic cancellation for Liouville and the standard summatory
estimate associated with a fixed zero-free half-plane. Related
multiplicative-function models appear in \cite{Round26aymone,Round26aymone-sidoravicius}.

\begin{figure}[tb]
\centering
\begin{tikzpicture}[>=Latex,x=8cm,y=1cm,font=\small]
\fill[blue!6] (.5,-.8) rectangle (1.15,.8);
\draw[->] (.40,0)--(1.2,0) node[right] {$\Re s$};
\draw[dashed] (.5,-1)--(.5,1) node[above] {$1/2$};
\draw[red!70!black,thick] (.76,-.7)--(.76,.7);
\fill[red!70!black] (.76,0) circle[x radius=.006,y radius=.045];
\node[below] at (.76,-.75) {real pole $\beta$};
\draw[blue!70!black,thick] (1,0) circle[x radius=.007,y radius=.055];
\node[above,align=center] at (1,.35) {simple zero of $F_f$\\double zero of $H_f$};
\node[below,align=center] at (.82,-1.65) {all logarithmic savings coexist with\\energy exponent $2\max(\Theta,\beta)$};
\end{tikzpicture}
\caption{Analytic structure of the deterministic multiplicative construction. The pole at the prescribed real point $\beta$ is present independently of the location of unknown zeta zeros.}
\label{round26:fig:rh-pole}
\end{figure}

\section{Rounding prime weights and computing the Euler product}

For each prime $p>y_0$, set $q_p=\tfrac12p^{\beta-1}$ and define
\[
 S(t)=\sum_{y_0<p\le t}q_p,
 \qquad I_p=\lfloor S(p)\rfloor-\lfloor S(p)-q_p\rfloor.
\]
Each $I_p$ is either zero or one, since $0<q_p<1$. Set $f(p)=-1$
at every prime at most $y_0$, and $f(p)=2I_p-1$ at larger primes.
Extend completely multiplicatively. Let $\mathcal S=\{p>y_0:I_p=1\}$
be the set of primes where the sign changes. Telescoping the floors
gives, for every real $t$,
\begin{equation}\label{round26:eq:discrepancy}
 E(t)=\sum_{y_0<p\le t}(I_p-q_p)=\lfloor S(t)\rfloor-S(t)\in(-1,0].
\end{equation}

Define the centered prime series $A(s)=\sum_{p>y_0}(I_p-q_p)p^{-s}$.
Summation by parts and \eqref{round26:eq:discrepancy} give
\begin{equation}\label{round26:eq:centered}
 A(s)=s\int_1^\infty E(t)t^{-s-1}\dd t\qquad(\R s>0).
\end{equation}
The integral converges locally uniformly in that half-plane and defines
a holomorphic function. For every real $\sigma>\beta$ it follows that
\begin{equation}\label{round26:eq:thin}
 \sum_{p>y_0}I_pp^{-\sigma}
   =\frac12\sum_{p>y_0}p^{-(\sigma+1-\beta)}+A(\sigma)<\infty.
\end{equation}
In particular $\sum_{p\in\mathcal S}1/p<\infty$, and the sum
$\sum_p(1-f(p)\lambda(p))/p$ equals $2\sum_{p\in\mathcal S}1/p$.

Consider the Euler product
\[
 G(s)=\prod_{p\in\mathcal S}\frac{1+p^{-s}}{1-p^{-s}}.
\]
It converges absolutely for $\R s>\beta$ by \eqref{round26:eq:thin}.
On $\R s>1$, comparison of local factors gives
$F_f(s)=\zeta(2s)G(s)/\zeta(s)$. To factor $G$, subtract the
logarithmic Euler series of $\zeta(s+1-\beta)$ from that of $G$.
The result on $\R s>1$ is
\begin{align}
 C_f(s)={}&2A(s)-\sum_{p\le y_0}p^{-(s+1-\beta)}\notag\\
  &+2\sum_{p>y_0}I_p
       \sum_{\substack{k\ge3\\ k\,\mathrm{odd}}}\frac{p^{-ks}}k
       -\sum_p\sum_{k\ge2}\frac{p^{-k(s+1-\beta)}}k.
       \label{round26:eq:factorlog}
\end{align}
The first term is holomorphic by \eqref{round26:eq:centered}. The last two
sums converge absolutely and locally uniformly for $\R s>1/2$.
For a lower bound $\sigma>1/2$ on the real part, their sums of
absolute values are bounded by constants times $\sum_p p^{-3\sigma}$
and $\sum_p p^{-2(\sigma+1-\beta)}$, respectively. Both exponents
exceed one. Thus $R_f=\exp C_f$ is holomorphic and nonvanishing on
$\R s>1/2$, and is positive at real points there. Exponentiating
\eqref{round26:eq:factorlog} gives $G(s)=\zeta(s+1-\beta)R_f(s)$ on
$\R s>1$, proving \eqref{round26:eq:factorization} by continuation.

The simple pole of $\zeta$ at~$1$ now gives
\begin{equation}\label{round26:eq:abel}
 \frac{F_f(s)}{\zeta(s)c_{y_0}(s)}
 =\frac{\zeta(2)\zeta(2-\beta)R_f(1)}{c_{y_0}(1)}(s-1)^2
       +O((s-1)^3).
\end{equation}
The coefficient is positive, and the zero of $F_f$ itself at~$1$
is simple. At $s=\beta$, one has $\zeta(\beta)<0$, by the alternating
series formula on $0<\beta<1$. The shifted zeta factor contributes
a genuine pole with residue
\begin{equation}\label{round26:eq:residue}
 \operatorname*{Res}_{s=\beta}F_f(s)
    =\frac{\zeta(2\beta)R_f(\beta)}{\zeta(\beta)}\ne0.
\end{equation}
The other denominator in \eqref{round26:eq:abel} is holomorphic and nonzero
at $\beta$, so that quotient also has a simple pole there.

The local factors of $H_f$ at $p\le y_0$ are one. At a larger prime in
$\mathcal S$ the factor is also one. At a larger prime outside
$\mathcal S$ it is
\[
 \frac{1-p^{-s}}{1+p^{-s}}=1+2\sum_{j\ge1}(-1)^jp^{-js}.
\]
Thus the coefficients are $2^{\omega(n)}\lambda(n)$ on integers
whose prime factors exceed $y_0$ and are outside $\mathcal S$,
and zero elsewhere. Here $\omega(n)$ counts distinct prime factors.

\section{Summatory bounds and the exact energy exponent}

Let $h$ be the multiplicative function with $h(1)=1$ and with
$h(p^j)=2$ for $p\in\mathcal S$, $j\ge1$, and zero at other prime
powers. Its coefficients are nonnegative and its Dirichlet series
is $G$. Local factors give the exact convolution
\begin{equation}\label{round26:eq:convolution}
 f=\lambda*h.
\end{equation}
Choose $\sigma_0\in(\beta,1)$, so that $G(\sigma_0)<\infty$.
The classical Davenport estimate at additive frequency zero gives
$L_\lambda(x)\ll_A x/(\log(2x))^A$ for every $A>0$. The Liouville
form is stated in \cite[\S1, after Lemma~1.4]{Round26mrt}, with attribution
to \cite{Round26davenport}. In the convolution sum
$L_f(x)=\sum_{d\le x}h(d)L_\lambda(x/d)$, use the logarithmic estimate
for $d\le\sqrt x$ and the bound $|L_\lambda(x/d)|\le x/d$ for the
remaining terms. This gives
\begin{align*}
 |L_f(x)|&\le C_A G(1)\frac{x}{(\log x)^A}
                +x\sum_{d>\sqrt x}\frac{h(d)}d\\
         &\le C_A G(1)\frac{x}{(\log x)^A}
                +G(\sigma_0)x^{(1+\sigma_0)/2}.
\end{align*}
Since $(1+\sigma_0)/2<1$, the second term is absorbed in the first
for each fixed $A$, with a constant valid for all $x\ge2$.
The same estimate holds uniformly for the additive twist
$\sum_{n\le x}f(n)e^{2\pi i n\alpha}$. In the convolution sum
the inner frequency becomes $d\alpha$, so Davenport's uniform
estimate applies with the same constant.
Lemma~\ref{round26:lem:prefix}, separating prefixes below $\sqrt N$ from
larger ones, also gives $Q_f(N)\ll_{A,f}N^2/(\log N)^A$ for every $A$.

Suppose $Q_f(N)=O(N^{2\alpha})$ for some $\alpha<\beta$.
Equation~\eqref{round26:eq:prefix-ineq} then gives $L_f(x)=O(x^\alpha)$,
and partial summation makes
\[
 s\int_1^\infty L_f(x)x^{-s-1}\dd x
\]
holomorphic for $\R s>\alpha$. It equals $F_f(s)$ initially on
$\R s>1$. The pole \eqref{round26:eq:residue} rules out this bound.
The same argument applies with any pole $\rho$ of $F_f$ to the
right of $\alpha$. If the energy limit superior in \eqref{round26:eq:exponent}
were below $2\R\rho$, choose $\alpha$ between half of that limit
superior and $\R\rho$. This would give the forbidden power bound.
Every zero $\rho$ of $\zeta$ with
$\R\rho>\beta$ is such a pole. Indeed both $2\rho$ and
$\rho+1-\beta$ have real part greater than one, so neither
zeta factor in the numerator of \eqref{round26:eq:factorization} vanishes.
Taking real parts of these zeros arbitrarily close to $\Theta$
when $\Theta>\beta$ gives the lower exponent $2\max(\Theta,\beta)$.

For the upper exponent, use the standard consequence of a fixed
zero-free half-plane that $L_\lambda(x)=O_\mu(x^\mu)$ for every
$\mu>\Theta$. The reciprocal-zeta contour argument is classical
\cite[Ch.~XIV]{titchmarsh}. One may pass from its M\"obius form to
Liouville using $\lambda(n)=\sum_{d^2\mid n}\mu_{\mathrm{Mob}}(n/d^2)$
and the convergence of $\sum_d d^{-2\mu}$ for $\mu>1/2$.
For $\mu>\max(\Theta,\beta)$, \eqref{round26:eq:convolution} therefore gives
\[
 |L_f(x)|\le C_\mu x^\mu\sum_{d\le x}\frac{h(d)}{d^\mu}
          \le C_\mu x^\mu G(\mu).
\]
Apply Lemma~\ref{round26:lem:prefix} and let $\mu$ decrease to
$\max(\Theta,\beta)$. This proves \eqref{round26:eq:exponent} and finishes
Theorem~\ref{round26:thm:construction}.

The same pole forbids a square-root bound for the exponential sum
$B_f(X)=\sum_{n\ge1}f(n)e^{-n/X}$. On $\R s>1$, absolute convergence
and Euler's gamma integral give
\[
 \int_0^\infty B_f(X)X^{-s}\frac{\dd X}{X}=\Gamma(s)F_f(s).
\]
If $B_f(X)=O(X^{1/2+\varepsilon})$ with $1/2+\varepsilon<\beta$,
the left side would be holomorphic to the right of
$1/2+\varepsilon$, since $B_f(X)$ decays exponentially as $X\downarrow0$.
The nonzero gamma factor at $\beta$ preserves the pole on the right.

\section{The Abel zero and its leading coefficient}

\begin{proposition}\label{round26:prop:slope}
Let $g$ be a completely multiplicative sign function whose Dirichlet
series is holomorphic near~$1$ and satisfies $F_g(1)=0$. Then the zero
is simple, $F_g'(1)\ge\zeta(2)$, and equality holds exactly for $g=\lambda$.
Consequently the quotient $F_g/(\zeta c_{y_0})$ has a double zero at~$1$.
Its leading coefficient equals $\zeta(2)/c_{y_0}(1)$ exactly when
$g=\lambda$.
\end{proposition}
\begin{proof}
For real $\sigma>1$, the Euler products give
\[
 \frac{F_g(\sigma)}{F_\lambda(\sigma)}
   =\prod_{p:g(p)=1}\frac{1+p^{-\sigma}}{1-p^{-\sigma}}\ge1.
\]
Since $F_\lambda(1+\delta)\sim\zeta(2)\delta$, a zero of order
greater than one would contradict this inequality. Dividing by
$\delta$ and letting $\delta\downarrow0$ proves the derivative bound.
If $g(p)=1$ at any prime $p$, the ratio has limit inferior at least
$(p+1)/(p-1)>1$. Thus equality of derivatives forces every prime sign
to be $-1$, and complete multiplicativity gives $g=\lambda$.
The quotient assertion follows by dividing by the pole of $\zeta$
and by the nonzero value $c_{y_0}(1)$.
\end{proof}

The constructed function attains the maximal zero order, while its
slope is $\zeta(2)G(1)>\zeta(2)$. The strict inequality follows from
the presence of a changed prime. In fact $S(t)\to\infty$, since the
sum of reciprocal primes diverges and $p^{\beta-1}\ge p^{-1}$.
Thus complete multiplicativity, agreement at any prescribed finite
prime prefix, every logarithmic saving, and maximal Abel vanishing
together permit energy growth above $N^{1+\varepsilon}$. The exact
Abel coefficient singles out Liouville within the sign class. The
square-root upper bound for Liouville itself remains open.
